\section{2023 Geometry \& Topology}

\subsection{Exam}

\begin{exercise}
\label{geometry:2023:1}
(a) Let $G$ be a Lie group, $\mathfrak { g }$ be its Lie algebra. The Maurer-Cartan form on $G$ is the unique left-invariant $\mathfrak { g }$ -valued 1-form such that $\omega | _ { e } : T _ { e } G \to \mathfrak { g }$ is the identity map, where $e$ is the identity in $G$ . Show that the Maurer-Cartan form $\omega$ satisfies the Maurer-Cartan equation $d \omega + \frac { 1 } { 2 } [ \omega , \omega ] = 0$.

(b) Let $G$ be a matrix group $G L ( n , \mathbb { R } )$ , give detailed computation to find the Maurer-Cartan form $\omega$ .

(c) Let $S E ( 3 )$ be the special Euclidean group, i.e. it contains all transformations of $\mathbb { R } ^ { 3 }$ (as 3-dim Euclidean space) of the form ${ \bf x } \mapsto t + A { \bf x }$ , where $\textbf { t } \in \mathbb { R } ^ { 3 }$ and $A \in S O ( 3 )$ . Find an expression of the Maurer Cartan form $\omega$ of $S E ( 3 )$ , also check it satisfies the standard Cartan structure equations in Euclidean space.
\end{exercise}

\begin{solution}
\textbf{Prerequisites:}
We use the left Maurer--Cartan form convention. If $X\in\mathfrak{g}$, denote by
$X^L$ the corresponding left-invariant vector field on $G$. For a
$\mathfrak{g}$-valued 1-form $\omega$, the bracket is the
$\mathfrak{g}$-valued 2-form
\[
    [\omega,\omega](U,V)
    =
    [\omega(U),\omega(V)]-[\omega(V),\omega(U)]
    =
    2[\omega(U),\omega(V)].
\]

\textbf{Part (a).}
By definition of the Maurer--Cartan form,
\[
    \omega(X^L)=X
\]
for every $X\in\mathfrak{g}$. Hence $\omega(X^L)$ and $\omega(Y^L)$ are constant
$\mathfrak{g}$-valued functions. Using the formula for the exterior derivative of
a 1-form,
\[
    d\omega(X^L,Y^L)
    =
    X^L(\omega(Y^L))-Y^L(\omega(X^L))-\omega([X^L,Y^L]).
\]
The first two terms vanish, and $[X^L,Y^L]=[X,Y]^L$. Therefore
\[
    d\omega(X^L,Y^L)=-[X,Y].
\]
On the other hand,
\[
    \frac{1}{2}[\omega,\omega](X^L,Y^L)=[X,Y].
\]
Thus
\[
    \left(d\omega+\frac{1}{2}[\omega,\omega]\right)(X^L,Y^L)=0.
\]
Since left-invariant vector fields span every tangent space and both sides are
left-invariant tensorial 2-forms, the Maurer--Cartan equation holds on all of
$G$.

\textbf{Part (b).}
For $G=GL(n,\mathbb{R})$, the tangent space at a matrix $A$ is naturally
$T_A G\simeq M_n(\mathbb{R})$. Left translation by $A^{-1}$ sends $A$ to the
identity and has differential
\[
    d(L_{A^{-1}})_A(V)=A^{-1}V.
\]
Therefore the left Maurer--Cartan form is
\[
    \omega_A(V)=A^{-1}V.
\]
In matrix-valued differential form notation this is
\[
    \omega=A^{-1}dA.
\]
To check the equation directly, differentiate $AA^{-1}=I$ to get
\[
    d(A^{-1})=-A^{-1}(dA)A^{-1}.
\]
Hence
\[
    d\omega
    =
    d(A^{-1})\wedge dA
    =
    -A^{-1}dA\wedge A^{-1}dA
    =
    -\omega\wedge\omega.
\]
For matrix Lie algebras, $\frac12[\omega,\omega]=\omega\wedge\omega$, so
$d\omega+\frac12[\omega,\omega]=0$.

\textbf{Part (c).}
Write an element of $SE(3)$ as a block matrix
\[
    g=
    \begin{pmatrix}
        A & t\\
        0 & 1
    \end{pmatrix},
    \qquad A\in SO(3),\quad t\in\mathbb{R}^3.
\]
Then
\[
    g^{-1}=
    \begin{pmatrix}
        A^T & -A^Tt\\
        0 & 1
    \end{pmatrix},
    \qquad
    dg=
    \begin{pmatrix}
        dA & dt\\
        0 & 0
    \end{pmatrix}.
\]
Thus the Maurer--Cartan form is
\[
    \omega=g^{-1}dg
    =
    \begin{pmatrix}
        A^TdA & A^Tdt\\
        0 & 0
    \end{pmatrix}.
\]
Set
\[
    \Omega=A^TdA\in\Omega^1(SE(3);\mathfrak{so}(3)),
    \qquad
    \theta=A^Tdt\in\Omega^1(SE(3);\mathbb{R}^3).
\]
Then
\[
    \omega=
    \begin{pmatrix}
        \Omega & \theta\\
        0 & 0
    \end{pmatrix}.
\]
The matrix Maurer--Cartan equation $d\omega+\omega\wedge\omega=0$ gives the two
block equations
\[
    d\Omega+\Omega\wedge\Omega=0,
    \qquad
    d\theta+\Omega\wedge\theta=0.
\]
Writing $\Omega=(\omega_{ij})$ and $\theta=(\theta_i)$, these become
\[
    d\theta_i+\sum_j\omega_{ij}\wedge\theta_j=0,
    \qquad
    d\omega_{ij}+\sum_k\omega_{ik}\wedge\omega_{kj}=0.
\]
These are precisely the Euclidean Cartan structure equations: zero torsion and
zero curvature for the standard moving frame in Euclidean space.
\end{solution}

\begin{exercise}
\label{geometry:2023:2}
Let $M_1$, $M_2$ be closed, connected, oriented 3-manifolds with the first fundamental groups
$$
\pi_ {1} (M _ {1}) = \mathbb {Z} _ {3} \oplus \mathbb {Z} ^ {2}, \quad \pi_ {1} (M _ {2}) = \mathbb {Z} _ {6} \oplus \mathbb {Z} ^ {3},
$$
(a) Find all homology groups $H _ { n } ( M _ { 1 } , \mathbb { Z } )$ and $H _ { n } ( M _ { 2 } , \mathbb { Z } )$ .

(b) Find all homology groups $H _ { n } ( M _ { 1 } \times M _ { 2 } , \mathbb { Q } )$ .

(c) Does there exist a closed connected oriented 3-manifold $M$ with
$$
\pi_ {1} (M) = \mathbb {Z} _ {3} \oplus \mathbb {Z} ^ {2} \quad \mathrm {o r} \quad \pi_ {1} (M) = \mathbb {Z} _ {6} \oplus \mathbb {Z} ^ {3}?
$$
\end{exercise}

\begin{solution}
\textbf{Prerequisites:}
For a connected space, $H_0\cong\mathbb{Z}$. For a closed oriented 3-manifold,
$H_3\cong\mathbb{Z}$ and Poincare duality gives $H_2\cong H^1$. The universal
coefficient theorem gives
\[
    H^1(M;\mathbb{Z})\cong\operatorname{Hom}(H_1(M;\mathbb{Z}),\mathbb{Z}).
\]
Also $H_1(M;\mathbb{Z})$ is the abelianization of $\pi_1(M)$.

\textbf{Part (a).}
The two given fundamental groups are already abelian, so
\[
    H_1(M_1;\mathbb{Z})\cong \mathbb{Z}_3\oplus\mathbb{Z}^2,
    \qquad
    H_1(M_2;\mathbb{Z})\cong \mathbb{Z}_6\oplus\mathbb{Z}^3.
\]
For $M_1$,
\[
    H_2(M_1;\mathbb{Z})
    \cong
    H^1(M_1;\mathbb{Z})
    \cong
    \operatorname{Hom}(\mathbb{Z}_3\oplus\mathbb{Z}^2,\mathbb{Z})
    \cong
    \mathbb{Z}^2.
\]
Therefore
\[
    H_n(M_1;\mathbb{Z})\cong
    \begin{cases}
        \mathbb{Z}, & n=0,3,\\
        \mathbb{Z}_3\oplus\mathbb{Z}^2, & n=1,\\
        \mathbb{Z}^2, & n=2,\\
        0, & \text{otherwise}.
    \end{cases}
\]
Similarly,
\[
    H_2(M_2;\mathbb{Z})
    \cong
    \operatorname{Hom}(\mathbb{Z}_6\oplus\mathbb{Z}^3,\mathbb{Z})
    \cong
    \mathbb{Z}^3,
\]
so
\[
    H_n(M_2;\mathbb{Z})\cong
    \begin{cases}
        \mathbb{Z}, & n=0,3,\\
        \mathbb{Z}_6\oplus\mathbb{Z}^3, & n=1,\\
        \mathbb{Z}^3, & n=2,\\
        0, & \text{otherwise}.
    \end{cases}
\]

\textbf{Part (b).}
Over $\mathbb{Q}$ the torsion disappears. Hence the Poincare polynomials of
$M_1$ and $M_2$ over $\mathbb{Q}$ are
\[
    P_{M_1}(t)=1+2t+2t^2+t^3,
    \qquad
    P_{M_2}(t)=1+3t+3t^2+t^3.
\]
The Kunneth theorem over the field $\mathbb{Q}$ gives
\[
    P_{M_1\times M_2}(t)
    =
    (1+2t+2t^2+t^3)(1+3t+3t^2+t^3).
\]
Thus
\[
    P_{M_1\times M_2}(t)
    =
    1+5t+11t^2+14t^3+11t^4+5t^5+t^6.
\]
Equivalently,
\[
    H_n(M_1\times M_2;\mathbb{Q})\cong
    \begin{cases}
        \mathbb{Q}, & n=0,6,\\
        \mathbb{Q}^5, & n=1,5,\\
        \mathbb{Q}^{11}, & n=2,4,\\
        \mathbb{Q}^{14}, & n=3,\\
        0, & \text{otherwise}.
    \end{cases}
\]

\textbf{Part (c).}
No such closed connected oriented 3-manifold exists for either group.

Here is the standard 3-manifold group obstruction. By the prime decomposition
theorem, the fundamental group of a closed oriented 3-manifold is a free product
of the fundamental groups of its prime factors. Since the proposed groups are
abelian, this free product can have at most one nontrivial factor. Thus the
manifold is prime up to summing with simply connected factors.

For a prime closed oriented 3-manifold, either the fundamental group is finite,
or it is torsion-free, except for the $S^1\times S^2$ case, whose fundamental
group is $\mathbb{Z}$. The proposed groups are infinite but contain nontrivial
torsion, namely $\mathbb{Z}_3$ or $\mathbb{Z}_6$. This is impossible. Therefore
neither $\mathbb{Z}_3\oplus\mathbb{Z}^2$ nor
$\mathbb{Z}_6\oplus\mathbb{Z}^3$ can occur as the fundamental group of a closed
connected oriented 3-manifold.
\end{solution}

\begin{exercise}
\label{geometry:2023:3}
(a) Let $f$ be a diffeomorphism group of a circle $S ^ { 1 }$ , assume $f$ has no fixed point and it is generated by a smooth vector field, show that $f$ must be conjugate to a rotation.

(b) Show that there is a diffeomorphism $f : S ^ { 1 } \to S ^ { 1 }$ , such that $f$ can not be generated by a smooth vector field but it is arbitrarily closed the identity map $i : S ^ { 1 } \to S ^ { 1 }$ in $C ^ { \infty }$ -topology.
\end{exercise}

\begin{solution}
\textbf{Part (a).}
Interpret ``generated by a smooth vector field'' as saying that $f$ is the time
one map of the flow of a smooth vector field $X$ on $S^1$. Since $f$ has no fixed
point, $X$ cannot vanish. Indeed, if $X(p)=0$, then the flow fixes $p$ for all
time, so $f(p)=p$.

Choose an angular coordinate $\theta$ on $S^1=\mathbb{R}/2\pi\mathbb{Z}$ and
write
\[
    X=v(\theta)\frac{\partial}{\partial\theta}.
\]
Because $X$ has no zero, after possibly reversing orientation we may assume
$v(\theta)>0$ everywhere. Define a new coordinate
\[
    s(\theta)
    =
    c\int_0^\theta \frac{du}{v(u)},
\]
where the positive constant $c$ is chosen so that $s(\theta+2\pi)=s(\theta)+2\pi$.
In the $s$-coordinate,
\[
    X=c\frac{\partial}{\partial s}.
\]
Therefore its time one map is
\[
    s\mapsto s+c \pmod{2\pi},
\]
which is a rotation. The coordinate change $\theta\mapsto s(\theta)$ is a smooth
diffeomorphism of the circle, so $f$ is smoothly conjugate to a rotation.

\textbf{Part (b).}
We construct diffeomorphisms arbitrarily $C^\infty$-close to the identity that
are not time one maps of autonomous smooth vector fields.

Fix a large integer $q$. Work with the lift of the circle
$S^1=\mathbb{R}/\mathbb{Z}$. Choose a smooth diffeomorphism
\[
    h:\left[0,\frac1q\right]\longrightarrow \left[0,\frac1q\right]
\]
which fixes the endpoints, agrees with the identity to infinite order at the
endpoints, is $C^\infty$-close to the identity, but is not the identity. Define a
smooth increasing lift $F:\mathbb{R}\to\mathbb{R}$ by first setting, on
$[0,1]$,
\[
    F(x)=x+\frac1q
    \quad
    \text{for }0\leq x\leq \frac{q-1}{q},
\]
and
\[
    F(x)=1+h\left(x-\frac{q-1}{q}\right)
    \quad
    \text{for }\frac{q-1}{q}\leq x\leq 1.
\]
Then extend by $F(x+1)=F(x)+1$. The infinite-order matching of $h$ with the
identity at the endpoints makes $F$ smooth. It defines an orientation-preserving
circle diffeomorphism $f$.

For $q$ large and $h$ sufficiently close to the identity, the displacement
$F(x)-x$ is positive and uniformly small, so $f$ has no fixed point and is
arbitrarily close to the identity in the $C^\infty$ topology. The endpoints of
the intervals form a periodic orbit of period $q$. However, on
$[0,1/q]$,
\[
    F^q(x)=1+h(x),
\]
so $f^q$ is not the identity because $h\neq\operatorname{Id}$.

If $f$ were generated by a smooth vector field, Part (a) would imply that $f$ is
conjugate to a rotation. Since $f$ has a point of period $q$, that rotation would
have rational angle $1/q$ up to relabeling, and then every point would be
$q$-periodic. This contradicts $f^q|_{I_0}=h\neq\operatorname{Id}$. Hence this
$f$ is not generated by any smooth vector field.
\end{solution}

\begin{exercise}
\label{geometry:2023:4}
(a) State the Leray-Hirsch theorem.

(b) Let
$$
\begin{array}{l} F l _ {k} (\mathbb {C} ^ {n}) = \{\mathrm {a l l\ } k \mathrm {-flags\ in\ } \mathbb {C} ^ {n} \} \\ = \{(F _ {0}, \dots , F _ {k}) | F _ {i} \mathrm {\ is\ an\ } i \mathrm {-dim\ subspace\ of\ } \mathbb {C} ^ {n}, \mathrm {\ s.t.\ } F _ {0} \subset F _ {1} \subset \dots \subset F _ {k} \subset \mathbb {C} ^ {n} \} \\ \end{array}
$$
Let $\Phi : F l _ { k } ( \mathbb { C } ^ { n } ) \to F l _ { k - 1 } ( \mathbb { C } ^ { n } )$ be the projection map sending a $k$ -flag $( F _ { 0 } , \dots , F _ { k } )$ to a $( k - 1 )$ -flag $( F _ { 0 } , \dots , F _ { k - 1 } )$ , it is known this is a fiber bundle. What is the fiber of $\Phi$ ?

(c) Compute the Euler Characteristic $\chi ( F l _ { n } ( \mathbb { C } ^ { n } ) )$ .
\end{exercise}

\begin{solution}
\textbf{Part (a).}
One standard form of the Leray--Hirsch theorem is the following. Let
\[
    F\hookrightarrow E\xrightarrow{\pi}B
\]
be a fiber bundle, and work with coefficients in a field, say $\mathbb{Q}$, or
more generally in a ring for which the relevant cohomology groups are free. Suppose
there are classes $e_1,\ldots,e_r\in H^*(E)$ whose restrictions to each fiber
$F_b=\pi^{-1}(b)$ form a basis of $H^*(F_b)$. Then the map
\[
    H^*(B)\otimes \operatorname{span}\{e_1,\ldots,e_r\}
    \longrightarrow
    H^*(E),
    \qquad
    \alpha\otimes e_i\mapsto \pi^*\alpha\smile e_i
\]
is an isomorphism of $H^*(B)$-modules. In particular, $H^*(E)$ is a free
$H^*(B)$-module with basis represented by the classes $e_i$.

\textbf{Part (b).}
Fix a $(k-1)$-flag
\[
    F_0\subset F_1\subset\cdots\subset F_{k-1}\subset\mathbb{C}^n.
\]
The fiber of $\Phi$ over this flag consists of all $k$-dimensional subspaces
$F_k$ satisfying
\[
    F_{k-1}\subset F_k\subset\mathbb{C}^n.
\]
Choosing such an $F_k$ is equivalent to choosing a 1-dimensional subspace of the
quotient vector space $\mathbb{C}^n/F_{k-1}$. Since
\[
    \dim_{\mathbb{C}}(\mathbb{C}^n/F_{k-1})=n-k+1,
\]
the fiber is
\[
    \mathbb{P}(\mathbb{C}^n/F_{k-1})
    \cong
    \mathbb{CP}^{\,n-k}.
\]

\textbf{Part (c).}
The full flag variety $Fl_n(\mathbb{C}^n)$ can be built by the successive
fibrations
\[
    Fl_k(\mathbb{C}^n)\longrightarrow Fl_{k-1}(\mathbb{C}^n),
    \qquad
    \text{fiber }\mathbb{CP}^{\,n-k}.
\]
Euler characteristic is multiplicative for such fiber bundles of finite CW
complexes. Also
\[
    \chi(\mathbb{CP}^m)=m+1.
\]
Therefore
\[
    \chi(Fl_n(\mathbb{C}^n))
    =
    \prod_{k=1}^{n}\chi(\mathbb{CP}^{\,n-k})
    =
    \prod_{k=1}^{n}(n-k+1)
    =
    n!.
\]
Equivalently, the Schubert cell decomposition has one cell for each permutation
in $S_n$, and all cells are even-dimensional, so the same answer is $|S_n|=n!$.
\end{solution}

\begin{exercise}
\label{geometry:2023:5}
On the Euclidean space $\mathbb { R } ^ { n }$ , we consider an $n - 1$ form $\alpha$ , which is of class $C ^ { 1 }$ , such that both $\alpha$ and $d \alpha$ are in $L ^ { 1 }$ . Show that $\int _ { \mathbb { R } ^ { n } } d \alpha = 0$ .
\end{exercise}

\begin{solution}
The point is to justify Stokes' theorem on the noncompact space by cutting off at
infinity.

Choose a smooth cutoff function $\chi_R:\mathbb{R}^n\to[0,1]$ such that
\[
    \chi_R=1 \text{ on } B_R,\qquad
    \chi_R=0 \text{ outside } B_{2R},\qquad
    |d\chi_R|\leq \frac{C}{R}.
\]
Since $\chi_R\alpha$ is compactly supported, Stokes' theorem gives
\[
    0
    =
    \int_{\mathbb{R}^n} d(\chi_R\alpha)
    =
    \int_{\mathbb{R}^n} d\chi_R\wedge\alpha
    +
    \int_{\mathbb{R}^n}\chi_R\,d\alpha.
\]
Hence
\[
    \left|\int_{\mathbb{R}^n}\chi_R\,d\alpha\right|
    =
    \left|\int_{\mathbb{R}^n} d\chi_R\wedge\alpha\right|
    \leq
    \frac{C'}{R}\int_{B_{2R}\setminus B_R}|\alpha|.
\]
Because $\alpha\in L^1(\mathbb{R}^n)$, the tail integrals
$\int_{B_{2R}\setminus B_R}|\alpha|$ tend to $0$ as $R\to\infty$. Therefore the
right-hand side tends to $0$.

On the other hand, $d\alpha\in L^1$ and $\chi_R\to 1$ pointwise, so dominated
convergence gives
\[
    \lim_{R\to\infty}\int_{\mathbb{R}^n}\chi_R\,d\alpha
    =
    \int_{\mathbb{R}^n}d\alpha.
\]
Combining the two limits yields
\[
    \int_{\mathbb{R}^n}d\alpha=0.
\]
\end{solution}

\begin{exercise}
\label{geometry:2023:6}
A complete Riemannian metric $g _ { i j }$ on a smooth manifold $M ^ { n }$ is called a gradient expanding Ricci soliton if there exists a smooth function $f$ on $M ^ { n }$ such that the Ricci tensor Ric of the metric $g$ is given by
$$
R i c + \operatorname {H e s s} f = \lambda g,
$$
for some negative constant $\lambda < 0$ . Show that if $M$ is compact, then a gradient expanding Ricci soliton must be an Einstein metric.
\end{exercise}

\begin{solution}
The soliton equation is
\[
    \operatorname{Ric}+\operatorname{Hess}f=\lambda g,
    \qquad \lambda<0.
\]
Taking the trace gives
\[
    R+\Delta f=n\lambda,
\]
where $R$ is the scalar curvature. Since $M$ is compact,
\[
    \int_M R\,dV=n\lambda\operatorname{Vol}(M).
\]
Thus the average value of $R$ is $n\lambda$.

We next use the standard scalar curvature identity for a gradient Ricci soliton:
\[
    \Delta_f R=2\lambda R-2|\operatorname{Ric}|^2,
    \qquad
    \Delta_f:=\Delta-\langle\nabla f,\nabla\cdot\rangle.
\]
Here is the derivation, to fix the signs. Taking the trace of the soliton
equation gives
\[
    R+\Delta f=n\lambda.
\]
Taking divergence of
$R_{ij}+\nabla_i\nabla_j f=\lambda g_{ij}$ and using the contracted Bianchi
identity gives
\[
    \frac12\nabla_jR+\nabla^i\nabla_i\nabla_j f=0.
\]
For a function,
\[
    \nabla^i\nabla_i\nabla_j f
    =
    \nabla_j\Delta f+\operatorname{Ric}_{jk}\nabla^k f.
\]
Since $\nabla_j\Delta f=-\nabla_jR$, we get
\[
    \nabla_jR=2\operatorname{Ric}_{jk}\nabla^k f.
\]
Taking divergence of this last identity yields
\[
    \Delta R
    =
    \langle\nabla R,\nabla f\rangle
    +2\operatorname{Ric}_{ij}\nabla_i\nabla_jf.
\]
Finally $\nabla_i\nabla_jf=\lambda g_{ij}-\operatorname{Ric}_{ij}$, so
\[
    \Delta R
    =
    \langle\nabla R,\nabla f\rangle
    +2\lambda R-2|\operatorname{Ric}|^2,
\]
which is exactly the displayed identity for $\Delta_fR$.

Let $p$ be a point where $R$ attains its minimum. At $p$ we have
\[
    \Delta_f R(p)=\Delta R(p)\geq 0.
\]
Therefore
\[
    0\leq 2\lambda R(p)-2|\operatorname{Ric}|^2(p).
\]
Using the Cauchy--Schwarz estimate
\[
    |\operatorname{Ric}|^2\geq \frac{R^2}{n},
\]
we get
\[
    0
    \leq
    2\lambda R(p)-2|\operatorname{Ric}|^2(p)
    \leq
    2\lambda R(p)-\frac{2}{n}R(p)^2.
\]
Since $\lambda<0$, this inequality forces
\[
    R(p)\geq n\lambda.
\]
Thus $R\geq n\lambda$ everywhere. But the average value of $R$ is exactly
$n\lambda$, so
\[
    R\equiv n\lambda.
\]
Substituting this constant value into the scalar soliton identity gives
\[
    0=2\lambda(n\lambda)-2|\operatorname{Ric}|^2,
\]
so
\[
    |\operatorname{Ric}|^2=n\lambda^2.
\]
Since $R=n\lambda$, equality holds in
$|\operatorname{Ric}|^2\geq R^2/n$. Equality in this estimate means
\[
    \operatorname{Ric}=\frac{R}{n}g=\lambda g.
\]
Hence $g$ is Einstein. The original soliton equation then gives
$\operatorname{Hess}f=0$; on a compact manifold this implies $f$ is constant up to
an additive constant on each connected component.
\end{solution}
